\chapter{まとめ}\chaplab{summary}
本報告書では，数理モデリングを用いて未来大学へ通学するバスの混雑緩和を目指す本プロジェクトの，前期における活動内容を報告した．

前期の活動は，数理モデリングの基礎固めから始まった．輪講会では『Pythonによる数値計算法の基礎』と『データ分析のための数理モデル入門』の2冊を教材として，数値計算法の基礎理論と，モデルの設計・パラメータ推定・評価という数理モデリングの一連の流れを学習し，メンバー全員が数理モデル構築のための共通の基礎を身につけた．並行してアイデア出しを行い，複数の候補テーマを比較・検討した結果，未来大学の学生にとって身近な問題である通学バスの混雑をテーマとして決定した．

テーマ決定後は，一部の停留所における乗降者数の計測やデータ分析を通じて混雑の実態を把握し，課題の難しさを整理したうえで，複数の解決策候補を比較検討した．その結果，新たな車両や運転手を必要としない，既存の便の一部を快速化し，短縮した時間を利用して混雑区間をもう一度運行する「快速化・2周運行」を解決策として提案するに至った．中間発表では，この課題設定と解決策のアイデアに対して聴講者から高い共感と関心を得た一方で，数理的手法の具体性や評価指標の明確化が今後の課題として指摘された．

後期は，モデル班とデータ班の2班体制で活動を進める．モデル班は現状のバス運行を再現する数理モデルの構築と評価指標の策定を，データ班は函館バス株式会社との連携によるデータの収集・整理を担い，シミュレーションによって快速化の効果を定量的に検証する．そのうえで，学生・一般利用者・バス事業者の三者への影響を踏まえた実現可能な運行案を提案し，未来大学の学生の快適な通学環境の実現に貢献することを目指す．
\bunseki{藤井優気}
